\section{Introducere}

\subsection{Motivație}

Algoritmi de segmentarea imaginilor au ca scop împărțirea unei imagini in segmente la nivelul pixelilor,
creând măști, ce folosesc pentru rezolvarea unor alte probleme, cum ar fi pentru crearea unui sistem de condus autonom
sau analiza imaginilor medicale.
Această împărțire poate fi semantică unde fiecare pixel primește clasa tipului de conținut pe care o reprezintă
sau în functie de obiect, unde fiecare obiect din imagine are atribuit o clasă, chiar dacă există mai multe obiecte de acelasi tip.


Diferența principală între detectarea de obiecte și segmentarea imaginilor la nivel de obiecte este generarea măstilor.
Astfel orice algoritm de segmentare a imaginilor la nivel de obiecte este și un algoritm de detectare a obiectelor, însă reciproca nu este adevarată


Astfel în această lucrare urmează să analizez mai multe metode de segmentare a imaginilor, in special cele la nivel de obiect, în scopul determinării valabilității algoritmului Slot Attention.

\subsection{Concepte abordate}

Această lucrare v-a evalua:

\begin{itemize}
  \item \textbf{Metode Clasice} K-means clustering, Mean-Shift Segmentation
  \item \textbf{Metode DeepLearning} U-Net, Slot Attention
\end{itemize}
